\chapter{Introduction}
\label{chap:introduction}

\section{Background}

Pretrained language models have changed the usual workflow of natural language processing.
Instead of training a separate architecture for every task, a general model is pretrained on a
large corpus and then adapted to a smaller downstream dataset. The Transformer architecture
\citep{vaswani2017attention} and models such as BERT \citep{devlin2019bert} established the
effectiveness of this transfer-learning paradigm. Later model families increased parameter counts
and broadened their capabilities, but the cost of updating and storing a complete copy of every
model also increased.

Full fine-tuning updates all parameters of a pretrained network. It remains conceptually simple,
but it requires optimizer states and gradients for the entire model and produces a complete
task-specific checkpoint. This is inconvenient when one base model must support many tasks,
domains, users, or deployment environments. Parameter-efficient fine-tuning (PEFT) instead freezes
most pretrained parameters and learns a small task-specific component. Adapter layers
\citep{houlsby2019parameter}, prefix tuning \citep{li2021prefix}, prompt tuning
\citep{lester2021power}, and low-rank updates are examples of this strategy.

LoRA \citep{hu2022lora} is especially attractive because it adds a low-rank update to selected
linear transformations. For a frozen weight matrix $W_0$, the adapted transformation is
\begin{equation}
  h = W_0 x + \Delta W x,
  \qquad
  \Delta W = BA,
  \label{eq:lora-basic}
\end{equation}
where the rank $r$ of $A$ and $B$ is much smaller than the dimensions of $W_0$. Only the low-rank
factors are trained. This design is easy to implement, creates small task-specific checkpoints,
and can be merged into the base weights for inference.

The rank $r$ controls both capacity and cost. A low rank may underfit a difficult task, while a
high rank increases parameters and can overfit small data. In common practice, however, the same
rank is assigned to every targeted matrix and is selected by convention or limited tuning.
This raises a direct question: should a dataset with 100 examples receive the same adapter capacity
as one with 10,000 examples?

\section{From LoRA to AdaLoRA}

AdaLoRA \citep{zhang2023adalora} addresses uniform allocation by representing updates in a
singular-value-like form and pruning less important components during training. Rather than giving
every matrix a permanent equal rank, AdaLoRA starts from a larger budget and progressively
allocates capacity according to importance. The method is theoretically appealing because model
layers and projections do not contribute equally to a downstream task.

Despite that motivation, fixed-rank LoRA remains easier to use. AdaLoRA requires an initial rank,
a target rank, a total training-step estimate, and a schedule describing when budget reduction
begins, how often allocation is updated, and when the final budget is fixed. Importance estimates
computed from short low-resource runs may also be noisy. In such settings the allocator has only a
small number of gradient observations before it must make pruning decisions.

It is difficult to prove a broad statement such as ``AdaLoRA is not used'' from the experiments in
this thesis. The defensible technical claim is narrower: AdaLoRA exposes more configuration choices
and runtime work than fixed-rank LoRA, while its benefit depends on the data and optimization
setting. Therefore, a useful extension should reduce arbitrary configuration and make the
allocation policy responsive to the downstream dataset.

\section{Motivation for Star-LoRA}

The starting observation of this project is that adapter capacity should reflect three broad
properties:
\begin{enumerate}[label=(\arabic*)]
  \item \textbf{Available evidence.} Larger datasets can support more trainable capacity, whereas
  very small datasets are vulnerable to overfitting and unstable estimates.
  \item \textbf{Distribution quality.} Class imbalance, rare tokens, and highly skewed token
  frequencies can make a nominal sample count overstate the effective amount of information.
  \item \textbf{Task complexity.} Diverse labels and varied input representations may require more
  adaptation capacity or a wider set of target transformations.
\end{enumerate}

\starlora\ operationalizes these ideas with measurable proxies. It analyzes the dataset before
training and maps the resulting profile to adapter settings. It also records a stability-aware
importance score during training to examine whether raw gradient-based importance is reliable.
The name is used in this thesis as shorthand for \emph{Stability- and Task-Aware Rank Adaptation}.

\section{Research Questions}

The empirical study is organized around five research questions.
\begin{description}[style=nextline,leftmargin=1.4cm]
  \item[RQ1] Does Star-LoRA improve standard AdaLoRA across different low-resource sample budgets?
  \item[RQ2] Does the policy select ranks that vary meaningfully with the amount and distribution
  of data?
  \item[RQ3] How does Star-LoRA compare with fixed-rank LoRA and full fine-tuning?
  \item[RQ4] What runtime and stability trade-offs accompany dataset-aware module expansion?
  \item[RQ5] Can an LLM-generated layer-wise policy provide a useful alternative to the
  deterministic heuristic?
\end{description}

\section{Contributions}

This undergraduate thesis makes the following contributions:
\begin{enumerate}
  \item It implements a lightweight dataset profiler for sample count, label distribution,
  token statistics, and embedding variation.
  \item It proposes a transparent bounded rank policy that combines dataset size, distribution,
  and complexity factors.
  \item It extends adapter selection from fixed attention projections to a dataset-dependent set
  of attention and MLP projections.
  \item It implements stability-aware logging based on an exponential moving average and
  variance normalization of weight--gradient importance.
  \item It evaluates full fine-tuning, LoRA, AdaLoRA, and Star-LoRA over five data budgets and
  multiple random seeds.
  \item It reports a negative exploratory result for LLM-assisted layer allocation and separates
  genuine LLM outputs from failed API calls that used heuristic fallback.
\end{enumerate}

\section{Thesis Organization}

\Cref{chap:background} reviews PEFT, LoRA, AdaLoRA, and dynamic rank methods.
\Cref{chap:motivation} develops the design argument for dataset-aware allocation.
\Cref{chap:method} describes Star-LoRA mathematically and algorithmically.
\Cref{chap:implementation} documents the prototype.
\Cref{chap:experiments} defines the experimental protocol.
\Cref{chap:results} reports the results.
\Cref{chap:discussion} discusses interpretation, limitations, and the exploratory LLM approach.
\Cref{chap:conclusion} concludes the thesis.
